\chapter{Designing Sound Static Analyses}
\label{ch:designing-sound-analyses}

A range query can return a perfectly plausible answer and still miscompile a program.
The transfer may have modeled unbounded addition while the instruction wraps, the answer may describe an earlier version of the function, or the client may have mistaken a possible value for a necessary one.
None of these defects is repaired by a faster worklist.
The engineering problem is to connect a statement about executions to a particular transformation of a particular IR revision.

Chapters~\ref{ch:ch09-dataflow-analysis}--\ref{ch:ch11-abstract-interpretation} supply the equation-solving machinery.
Here we build its specification: what executions are collected, what an abstract element denotes, which local obligations make the result sound, and what a consumer may conclude.
The useful deliverable is an analysis contract that can be reviewed independently of the scheduling heuristics in its implementation.

\section{Start with the client and the executions}
\label{sec:ch42-client-first}

Suppose a compiler wants to remove the zero-divisor check before an integer division.
The required proposition is universal: every reachable state at that instruction has a nonzero divisor.
A warning tool asks a different question when it reports that a divisor might be zero.
Its report may mean that the analysis failed to exclude zero; it need not establish an execution that divides by zero.
The two tools can share a range engine while exposing different answers.

First choose the concrete semantics.
For the initial examples, values are mathematical integers, execution is sequential, and assignments have no hidden effects.
A state $\sigma$ maps the variables in scope to integers.
An edge $e$ has a transition relation $\tau_e\subseteq\Sigma\times\Sigma$, where $\Sigma$ is the state space.
An assignment edge relates the old store to the store with one binding replaced.
A guard edge relates a store to itself only when the condition holds.
Nondeterministic input is a relation with several possible successors, not an arbitrary fixed constant.

An actual compiler extends the state with memory, object lifetime, and whatever control state affects the instruction semantics.
If exceptions or traps are observable, give them explicit outcomes or outgoing edges.
Do not model a trapping instruction merely by omitting its normal successor and then conclude that the instruction can be erased.
The absence of a normal continuation says nothing about the observability of the trap itself.

The entry description $I_B\subseteq\Sigma$ is nonempty only at permitted entry points.
It includes all inputs allowed by the function's contract, rather than inputs seen in a profile.
For ordinary functions there is one entry; exception handlers and coroutine resumptions can require additional boundaries.
The \term{collecting semantics} assigns to every point $B$ the set $C_B$ of states occurring there on some finite execution prefix.
With $\operatorname{post}_e(S)=\{\sigma'\mid \exists\sigma\in S:(\sigma,\sigma')\in\tau_e\}$, it satisfies
\begin{equation}
 C_B=I_B\cup\bigcup_{e=(P,B)}\operatorname{post}_e(C_P).
 \label{eq:ch42-collecting}
\end{equation}
The intended solution of Equation~\ref{eq:ch42-collecting} is the least one under inclusion.
It collects only states justified by finite prefixes from an entry, including prefixes of executions that never terminate.
This choice matters to compilers: a loop with no return can still perform observable stores.

A set of states is sufficient for many safety properties, but not for every compiler question.
The property that an instruction eventually executes concerns futures, and a bound on total running time needs more than a set of stores at one point.
One can collect traces, costs, or state pairs instead, but must change the specification accordingly.
The choice of concrete object comes before the choice of lattice~\cite{nielson-1999-principles}.

\section{Give every representation a meaning}
\label{sec:ch42-representation}

Throughout this chapter and Chapters~\ref{ch:relational-path-analysis} and~\ref{ch:interprocedural-heap-analysis}, smaller means more precise.
Concrete sets use inclusion.
Abstract elements use $a\sqsubseteq b$ when $a$ describes no more behavior than $b$; in the canonical domains used here this is equivalent to $\gamma(a)\subseteq\gamma(b)$.
Thus $\bot$ describes no states, $\top$ describes all states, and $a\join b$ covers both alternatives.
This is the concretization order, rather than the reversed flat-domain convention used in some examples of Chapter~\ref{ch:ch10-lattices-fixed-points}.

For one integer variable, an interval $[l,u]$ denotes all integers between its endpoints.
Bounds may be infinite, and the empty interval has one canonical representation, $\bot$.
For several variables, a product of intervals denotes the Cartesian product of their value sets.
It does not denote only the tuples that happened to produce the endpoints.
In particular, $x\in[2,5]$ and $y\in[2,5]$ permit $x=2,y=5$, even if the program always assigns $y=x$.

The abstraction $\alpha(S)$ takes the least interval product covering $S$: each component uses the infimum and supremum of that variable's values, with $\alpha(\emptyset)=\bot$.
The defining relation is
\begin{equation}
 \alpha(S)\sqsubseteq a\quad\Longleftrightarrow\quad S\subseteq\gamma(a).
 \label{eq:ch42-adjunction}
\end{equation}
Equation~\ref{eq:ch42-adjunction} is useful as an implementation guide.
It makes abstraction a cover, not a sample, and identifies the best answer expressible by the representation.
Some useful abstractions have only a concretization or lack a computable best abstraction.
They can still support sound analysis; they must establish coverage directly rather than pretending to implement an unavailable $\alpha$.

Representation details belong in this specification.
Are endpoints inclusive?
Does an absent variable mean unconstrained, uninitialized, or out of scope?
Does an empty map represent one unconstrained store or no stores?
For the interval product here, an absent binding means $\top$ for that variable, while a separate bottom tag means no stores.
Conflating those two cases can turn an unknown argument into an unreachable function.

SSA names simplify this bookkeeping but do not eliminate semantic obligations.
A $\phi$-node selects the operand belonging to the incoming edge; it does not independently select arbitrary operands at a block entry.
If several $\phi$-nodes are evaluated by a relational domain, their edge-specific assignments must happen in parallel before alternatives are joined.
Sequentially overwriting names can destroy an old value still needed by another assignment.

\section{Derive and check local transfers}
\label{sec:ch42-local-soundness}

An abstract edge transformer $F_e^\#$ is sound when
\begin{equation}
 \operatorname{post}_e(\gamma(a))\subseteq\gamma(F_e^\#(a))
 \quad\text{for every }a.
 \label{eq:ch42-local-cover}
\end{equation}
Figure~\ref{fig:ch42-sound-square} shows the two routes in Equation~\ref{eq:ch42-local-cover}.
When a best abstraction exists, $\alpha\circ\operatorname{post}_e\circ\gamma$ is the most precise sound transformer available in that domain.
An implementation can use a cheaper, larger result.
The abstract-interpretation framework separates this local approximation from the global fixed-point argument~\cite{cousot-1977-abstract}.

\begin{figure}
\centering
\begin{tikzpicture}[node distance=17mm and 24mm]
\node[basic block] (ci) {$\gamma(a)$};
\node[basic block,right=of ci] (co) {$\operatorname{post}_e(\gamma(a))$};
\node[basic block,below=of ci] (ai) {$a$};
\node[basic block,below=of co] (ao) {$F_e^\#(a)$};
\draw[cfg edge] (ci) -- node[above,font=\small] {concrete step} (co);
\draw[cfg edge] (ai) -- node[below,font=\small] {abstract step} (ao);
\draw[cfg edge] (ai) -- node[left,font=\small] {denotes} (ci);
\draw[cfg edge] (ao) -- node[right,font=\small,align=left] {denotes a cover\\of the result} (co);
\end{tikzpicture}
\caption{Local soundness requires the abstract route to cover every concrete successor; equality between the routes is unnecessary.}
\label{fig:ch42-sound-square}
\end{figure}


For mathematical addition, $[l_1,u_1]+[l_2,u_2]=[l_1+l_2,u_1+u_2]$ is sound because addition is monotone in each operand.
For the guard $x<8$ on integers, intersect the interval for $x$ with $[-\infty,7]$.
If the result is empty, return bottom for the whole state.
The number $7$ comes from discrete integer semantics; the same replacement is wrong over the reals.
Unsupported predicates may leave the state unchanged, which loses precision but preserves coverage.

Repeated operands deserve special treatment.
Generic interval multiplication treats its two operands independently, whereas $x*x$ uses the same value twice.
For $x\in[-4,6]$, independent multiplication gives $[-24,36]$.
A square transfer gives $[0,36]$: zero is the minimum when the input crosses zero, and the maximum is the larger squared endpoint.
Both answers are sound, but only the second exploits the equality of the operands.
Neither reconstructs a lost disjunction in the input.

Consider a branch that assigns $x=-4$ on its left arm and $x=6$ on its right arm, then computes $y=x+5$ and $z=x*x$.
Table~\ref{tab:ch42-join-loss} computes the exact sets and interval results.
The actual values of $z$ are $16$ and $36$, but the joined interval input also admits $x=0$.
Even the best interval square transfer must therefore admit $z=0$.
Deleting a zero check on $z$ requires a more precise input description, such as retaining both branch alternatives until after the square.

\begin{table}
\centering
\caption{A sound join loses a disjunction before the square is evaluated.}
\label{tab:ch42-join-loss}
\begin{tabular}{p{0.25\linewidth}p{0.29\linewidth}p{0.34\linewidth}}
\toprule
Point or operation & Exact value set & Interval result \\
\midrule
Left arm, $x$ & $\{-4\}$ & $[-4,-4]$ \\
Right arm, $x$ & $\{6\}$ & $[6,6]$ \\
Join, $x$ & $\{-4,6\}$ & $[-4,6]$ \\
$y=x+5$ & $\{1,11\}$ & $[1,11]$ \\
$z=x*x$, after join & $\{16,36\}$ & $[0,36]$ with square transfer \\
Square each arm, then join & $\{16,36\}$ & $[16,36]$ \\
\bottomrule
\end{tabular}
\end{table}

Monotonicity and soundness are distinct obligations.
A transformer that returns bottom for every input is monotone but unsound for an executable assignment.
A locally sound transformer can also be implemented with precision choices that violate monotonicity; the usual ascending solver proofs then need additional care.
The baseline interface requires both properties, and requires $F_e^\#(\bot)=\bot$ for these execution-edge transformers.
Tests should exercise bottom, infinities, singleton values, and boundaries, not only typical finite intervals.

\section{Validate the answer rather than trusting the iteration history}
\label{sec:ch42-certificates}

The solver may use widening, priority queues, cached transfer results, or parallel work.
The final safety obligation has a simpler form.
Let $b_B$ cover $I_B$, and let $A_B$ be the proposed invariant at point $B$.
It is enough that $b_B\sqsubseteq A_B$ and that every edge satisfies $F_e^\#(A_P)\sqsubseteq A_B$.
These conditions say that the proposed result includes entry states and is closed under every abstract transition.
Such an assignment is a \term{post-fixed invariant}, using the convention $F(A)\sqsubseteq A$ for the global operator.
Some literature calls this a pre-fixed point; the displayed inequality is the contract used here.

Algorithm~\ref{alg:ch42-validate} checks this certificate on one immutable IR revision.
It does not need to reproduce the solver's choices or prove that the answer is the least one.
Its trusted components are the boundary abstraction, edge enumeration, transfer functions, and inclusion test.
Reusing a buggy transfer in the validator cannot detect that transfer's unsoundness, so this check complements local transfer verification.

\begin{algorithm}
\caption{Validate a proposed invariant on an immutable IR revision.}
\label{alg:ch42-validate}
\begin{pseudocode}
\Require revision $r$, graph $G$, boundary covers $b$, candidate $A$
\Require sound edge transfers and a sound inclusion test
\Ensure accept only if $A$ covers all executions from $b$ in $r$
\If{$A.\mathit{revision}\neq r$}
  \State \Return reject
\EndIf
\ForAll{points $B$ in $G$}
  \If{$b_B\not\sqsubseteq A_B$}
    \State \Return reject
  \EndIf
\EndFor
\ForAll{edges $e=(P,B)$ in $G$}
  \State $t\gets F_e^\#(A_P)$
  \If{$t\not\sqsubseteq A_B$}
    \State \Return reject
  \EndIf
\EndFor
\State \Return accept
\end{pseudocode}
\end{algorithm}

The correctness argument is induction on execution-prefix length.
An initial state is covered by the boundary check.
If a state at $P$ is covered, local soundness covers every successor in $F_e^\#(A_P)$, and the edge check places that successor inside $A_B$.
Every reachable state has a finite prefix, so all reachable states are covered.
No termination assumption about the analyzed program is used.

For the branch example, proposing $x\in[0,6]$ at the join fails the left-edge check because $[-4,-4]\not\sqsubseteq[0,6]$.
Proposing $[-4,6]$ passes both incoming checks.
At the square's output, a proposal of $[16,36]$ fails when the transfer is applied to the joined input $[-4,6]$.
Although $[16,36]$ is a true property of this particular program, this certificate and domain cannot establish it from those intermediate invariants.
Validation can reject a true claim whose proof information was discarded.

Let $N$ count program points, $E$ edges, $C_{\leq}$ the cost of an inclusion check, and $C_e$ the cost of edge $e$'s transfer.
The validator costs $O((N+E)C_{\leq}+\sum_e C_e)$ after the invariants exist.
An interval product over $V$ variables uses $O(V)$ endpoint comparisons per whole-state inclusion, rather than one constant-time operation.
With $b$-bit arbitrary-precision endpoints, arithmetic and comparison costs must also be charged; unit-cost bounds apply only when endpoints fit the chosen representation.

\section{Machine instructions change the proof}
\label{sec:ch42-machine}

The same notation $x+1$ can describe several concrete operations.
An unsigned eight-bit addition maps $255$ to $0$.
The mathematical transfer $[250,255]+1=[251,256]$ therefore misses a reachable result and is unsound for that instruction.
The exact output set is $\{0,251,252,253,254,255\}$, whose ordinary unsigned interval hull is $[0,255]$.
A wrapped interval representation can retain the two segments, but needs its own order, join, and transfer definitions.

For a checked signed addition, the normal result and the overflow outcome are separate successors.
For an IR operation whose overflow creates a special invalid value, a range domain that tracks only ordinary integers is insufficient by itself.
It needs a companion validity component or a precondition proving that the invalid case cannot arise.
For a language with undefined signed overflow, the optimizer's reasoning must be attached to the language and IR rules, rather than transferred mechanically from the unsigned case.

Conversions are equally important.
An interval for a wide integer does not automatically remain true after truncation, and signed comparison is not interchangeable with unsigned comparison on the same bits.
Floating-point transfers must account for NaNs, infinities, signed zero, and the permitted rounding behavior before they can justify a rewrite.
These are changes to the concrete semantics, not small inaccuracies absorbed by widening.

There is a practical audit method: for each opcode, record its input precondition, ordinary output relation, exceptional outcomes, and non-value effects.
Then state which components of the abstract state model each item.
An unmodeled effect requires a conservative fallback at that operation.
For example, an unknown call can retain facts about immutable scalar SSA values while forgetting facts about reachable mutable memory.
Forgetting everything is safe, but forgetting the wrong subset is not.

\section{Results have lifetimes and dependencies}
\label{sec:ch42-invalidation}

A true property of an old graph need not hold in a new graph.
Suppose an edit changes the left-arm assignment from $-4$ to $-5$.
The old join range $[-4,6]$ now misses a concrete value, and the old conclusion that $y=x+5$ is nonzero is false on that arm.
Even a transformation known to preserve program behavior can change the names and program points to which a cached answer refers.

An analysis cache should identify the IR revision and every prerequisite used to compute an answer.
Figure~\ref{fig:ch42-dependencies} shows a dependency chain in which an alias summary affects range information through load invalidation.
The reverse edges answer the invalidation question: which cached results read this changed fact?
Instruction identity alone is insufficient when a callee summary, entry assumption, or target arithmetic mode changes.

\begin{figure}
\centering
\begin{tikzpicture}[node distance=10mm]
\node[basic block] (body) {callee body revision};
\node[basic block,below=of body] (effect) {callee effect summary};
\node[basic block,below=of effect] (load) {caller load facts};
\node[basic block,below=of load] (range) {caller range invariant};
\node[basic block,below=of range] (rewrite) {check-removal decision};
\node[basic block,right=15mm of load] (alias) {alias result};
\draw[cfg edge] (body) -- (effect);
\draw[cfg edge] (effect) -- (load);
\draw[cfg edge] (alias) -- (load);
\draw[cfg edge] (load) -- (range);
\draw[cfg edge] (range) -- (rewrite);
\end{tikzpicture}
\caption{An arrow runs from a dependency to its reader. A changed callee can invalidate a scalar range proof through the intervening memory facts.}
\label{fig:ch42-dependencies}
\end{figure}


Algorithm~\ref{alg:ch42-invalidate} gives a conservative invalidation policy.
The dependency graph contains an edge $u\to v$ when cached result $v$ was computed using $u$.
Dependencies include the analyzed bodies and their versioned semantic options, not just other analyses.
The graph must already include every dependency used by each cached result.
New dependencies are registered when results are recomputed.

\begin{algorithm}
\caption{Invalidate the transitive readers of changed dependencies.}
\label{alg:ch42-invalidate}
\begin{pseudocode}
\Require complete reader graph $R$, changed dependency keys $D$
\Ensure every cached result transitively depending on $D$ is invalid
\State $W\gets D$; $Seen\gets D$
\While{$W\neq\emptyset$}
  \State remove $u$ from $W$
  \State mark the cached result for $u$ invalid, if one exists
  \ForAll{$v\in R[u]$}
    \If{$v\notin Seen$}
      \State insert $v$ into $Seen$ and $W$
    \EndIf
  \EndFor
\EndWhile
\end{pseudocode}
\end{algorithm}

On the diagram, changing the callee body invalidates its effect summary, then the caller's load facts, then the range result, and finally the cached check-removal decision.
Changing an unrelated function reaches none of those nodes if the recorded dependencies are complete.
The traversal takes $O(K+L)$ time for $K$ visited keys and $L$ traversed reader edges, with constant-time marks and adjacency lists.
It uses $O(K)$ temporary queue and membership storage.
This bound covers invalidation, not subsequent recomputation.

Restarting an ascending analysis from its old solution deserves a separate proof.
Adding behaviors can require enlarging the result, while deleting behaviors can leave obsolete facts that a grow-only solver never retracts.
More seriously, an old answer may not be a sound cover at all after a transfer changes.
The baseline is to invalidate, initialize from the specified boundaries and bottom, and solve again.
Incremental systems can retain unaffected components, but must account for changed equations, removed dependencies, and newly reachable paths.

\section{Expose proof-strength answers to consumers}
\label{sec:ch42-query-contract}

A useful query interface distinguishes unreachable, proved, refuted by a validated witness, and unknown.
These names refer to a precise proposition at a precise program point under explicit assumptions.
If $A_B=\bot$ in a validated result, the point is unreachable under those assumptions.
If $\gamma(A_B)\subseteq P$, the property $P$ is proved.
If the abstraction merely intersects the complement of $P$, the answer is unknown unless a concrete counterexample is independently established.

For $P$ equal to $z\neq0$, Table~\ref{tab:ch42-join-loss} yields unknown after the merged square.
For $P$ equal to $y\neq0$, it yields proved because $[1,11]$ excludes zero.
A diagnostic may explain that the first proof failed at a branch merge.
It must not describe the abstract value zero as an observed failing input.
Conversely, suppressing a warning because of a user preference must not upgrade an unknown answer into proof permission for the optimizer.

Proof permission is narrower than transformation permission.
Proving a divisor nonzero removes one failure condition; it does not establish that an entire division is removable, duplicable, or movable across side effects.
An integer division may have another exceptional operand pair, and evaluating the divisor may itself be observable.
The rewrite's own contract must cover those facts.
Likewise, a range proof at a use does not permit speculating the producing load onto a path where its address is invalid.

Analysis answers should therefore carry the proposition, program-point identity, revision, and assumptions that made them valid.
A guarded specialization can use a stronger entry assumption only below a dominating guard that establishes it.
The general path must retain behavior for inputs outside that assumption.
This is the point where abstract interpretation becomes a compiler interface rather than merely a mathematical computation.

\section{Exercises}
\label{sec:ch42-exercises}

\begin{enumerate}
\item \label{ex:ch42-square}
For branches assigning $x=-3$ and $x=8$, compute the interval for $x*x$ when joining before the square and when joining afterward.
Use the specialized square transfer.
Which answer proves that the result is nonzero, and why is the weaker answer still sound?
\item \label{ex:ch42-wrap}
An unsigned eight-bit value has range $[253,255]$ and is incremented by two modulo $256$.
Compute the exact output set and its ordinary interval hull.
Give the concrete input refuting the mathematical-integer transfer.
\item \label{ex:ch42-certificate}
A graph has entry $x=2$, an edge adding three, and then an exit.
Check candidate intervals $[2,2]$ at entry and $[2,4]$ at exit with Algorithm~\ref{alg:ch42-validate}.
Give the failing inclusion and a smallest replacement exit interval.
\item \label{ex:ch42-dependencies}
Dependencies are $a\to b$, $b\to c$, $a\to d$, and $e\to d$.
Starting from a change to $b$, give the invalidated set.
Then explain why adding a new dependency $e\to c$ without recording it can make a later change to $e$ unsafe.
\end{enumerate}

\section{Further reading}
\label{sec:ch42-reading}

Cousot and Cousot provide the foundational connection between concrete semantics, abstraction, and sound fixed-point approximation~\cite{cousot-1977-abstract}.
Nielson, Nielson, and Hankin develop semantic specifications alongside analysis construction; read them to distinguish correctness of a solver from correctness of the equations it solves~\cite{nielson-1999-principles}.
Kildall's framework explains why many compiler analyses share a propagation engine while computing different properties~\cite{kildall-1973-dataflow}.
